\section{Extensions}
\label{sec:extensions}

\subsection{The Cole--Hopf transformation}
\label{subsec:cole-hopf}

This extension reintroduces viscosity and solves the viscous Burgers'
equation \eqref{eq:viscous-burgers} exactly, using the potential-flow
structure noted in Remark~\ref{rem:vorticity-free}.

\begin{theorem}[Cole--Hopf transformation]
\label{thm:cole-hopf}
Let $u$ solve $u_t+u\,u_x=\nu\,u_{xx}$, and define $\varphi>0$ by
\begin{equation}
  \label{eq:cole-hopf-def}
  u = -2\nu\,\frac{\varphi_x}{\varphi}, \qquad \text{i.e.} \qquad \varphi(x,t) = \exp\left(-\frac{1}{2\nu}\int_0^x u(y,t)\, dy\right).
\end{equation}
Then $\varphi$ solves the heat equation $\varphi_t = \nu\, \varphi_{xx}$.
\end{theorem}

\begin{proof}
Write $\psi = \ln\varphi$, so $u=-2\nu\,\psi_x$. Substituting into the
viscous Burgers' equation and dividing by $-2\nu$ gives
\begin{equation}
  \psi_{xt} - 2\nu\,\psi_x\psi_{xx} - \nu\,\psi_{xxx} = 0,
\end{equation}
which is exactly $\partial_x\big[\psi_t - \nu\,\psi_{xx} - \nu\,(\psi_x)^2\big]=0$.
Hence $\psi_t - \nu\,\psi_{xx} - \nu\,(\psi_x)^2 = C(t)$ for some function
$C$ of $t$ alone; since $u=-2\nu\,\psi_x$ is unchanged if $\psi$ is
shifted by a function of $t$, $\varphi$ may be rescaled by
$\exp\!\big(-\!\int_0^t C\big)$ without affecting $u$, so we may take
$C\equiv0$:
\begin{equation}
  \label{eq:psi-equation}
  \psi_t = \nu\,\psi_{xx} + \nu\,(\psi_x)^2.
\end{equation}
Since $\varphi=e^\psi$ gives $\varphi_t=\varphi\,\psi_t$ and
$\varphi_{xx}=\varphi\big(\psi_{xx}+(\psi_x)^2\big)$, equation
\eqref{eq:psi-equation} is equivalent to $\varphi_t=\nu\,\varphi_{xx}$.
\end{proof}

\subsection{Exact solution via the heat kernel}
\label{subsec:cole-hopf-solution}

By Theorem~\ref{thm:cole-hopf}, solving \eqref{eq:viscous-burgers} reduces
to solving the heat equation for $\varphi$ with initial data
\begin{equation}
  \label{eq:phi0}
  \varphi_0(x) = \exp\left(-\frac{1}{2\nu}\int_0^x u_0(y)\, dy\right),
\end{equation}
which is given by convolution with the heat kernel,
\begin{equation}
  \label{eq:phi-heat-kernel}
  \varphi(x,t) = \int_{\mathbb{R}} \frac{1}{\sqrt{4\pi\nu t}}\,
    \exp\left(-\frac{(x-y)^2}{4\nu t}\right) \varphi_0(y)\, dy.
\end{equation}
Substituting \eqref{eq:phi0}--\eqref{eq:phi-heat-kernel} into
$u=-2\nu\,\varphi_x/\varphi$ and differentiating under the integral gives
the explicit \emph{Cole--Hopf formula}
\begin{equation}
  \label{eq:cole-hopf-formula}
  u(x,t) = \frac{\displaystyle\int_{\mathbb{R}} \frac{x-y}{t}\, \exp\left(-\frac{G(x,y,t)}{2\nu}\right) dy}
                {\displaystyle\int_{\mathbb{R}} \exp\left(-\frac{G(x,y,t)}{2\nu}\right) dy},
  \qquad
  G(x,y,t) = \int_0^y u_0(s)\, ds + \frac{(x-y)^2}{2t}.
\end{equation}

\begin{corollary}[Global smoothness]
\label{cor:global-smoothness}
For every bounded measurable $u_0$ and every $\nu>0$, formula
\eqref{eq:cole-hopf-formula} defines the unique bounded solution of
\eqref{eq:viscous-burgers}, and this solution is $C^\infty$ for all
$t>0$.
\end{corollary}

Unlike the inviscid equation (Proposition~\ref{prop:shock-time}),
solutions of \eqref{eq:viscous-burgers} stay smooth for all time, for
any $\nu>0$ however small: advection and viscous diffusion together do
not by themselves produce singularities (Remark~\ref{rem:navier-stokes-relevance}).

\subsection{The vanishing viscosity limit}
\label{subsec:vanishing-viscosity}

The Cole--Hopf solution also explains why the Lax entropy condition of
Section~\ref{sec:entropy-condition} is the physically correct selection
principle, rather than an arbitrary one.

\begin{theorem}[Vanishing viscosity limit; Hopf, 1950]
\label{thm:vanishing-viscosity}
Let $u_0\in L^\infty(\mathbb{R})$, let $u^\nu$ denote the solution of
\eqref{eq:viscous-burgers} with data $u_0$ given by
\eqref{eq:cole-hopf-formula}, and let $u$ denote the entropy solution of
\eqref{eq:inviscid-burgers} with the same data
(Theorem~\ref{thm:riemann-solution}, extended to general $u_0$ by the same
selection principle). Then $u^\nu \to u$ pointwise almost everywhere,
and in $L^1_{\mathrm{loc}}$, as $\nu \to 0^+$ \cite{hopf1950}.
\end{theorem}

The proof (a Laplace-method asymptotic analysis of
\eqref{eq:cole-hopf-formula} as $\nu\to0$, showing the integral
concentrates near the $y$ minimising $G(x,y,t)$) is not reproduced here.
Its significance is conceptual: the entropy condition of
Definition~\ref{def:lax-entropy} is not an ad hoc rule imposed to force
uniqueness, but exactly the criterion selected by the physically natural
limit of vanishing viscosity --- the same limit by which the inviscid
equation \eqref{eq:inviscid-burgers} was justified as an idealisation of
\eqref{eq:viscous-burgers} in Section~\ref{sec:motivation}.

\subsection{General convex flux}
\label{subsec:general-flux}

The theory of Sections~\ref{sec:weak-solutions}--\ref{sec:riemann-problem}
was developed for the specific flux $f(u)=u^2/2$, but, as noted in
Section~\ref{subsec:conservation-law-framework}, none of it relies on
this choice beyond convexity of $f$. This extension records how the main
results generalise to an arbitrary $f \in C^2(\mathbb{R})$ with
$f''>0$ everywhere.

\begin{proposition}[General Rankine--Hugoniot and Lax conditions]
\label{prop:general-rh-lax}
For a weak solution of $u_t+f(u)_x=0$ with a jump from $u_\ell$ to $u_r$
along $x=s(t)$, the Rankine--Hugoniot condition reads
\begin{equation}
  \label{eq:general-rh}
  s'(t) = \frac{f(u_\ell)-f(u_r)}{u_\ell-u_r},
\end{equation}
and the Lax entropy condition reads $f'(u_\ell) > s'(t) > f'(u_r)$. Since
$f$ is convex, $f'$ is strictly increasing, so this is equivalent to
$u_\ell > u_r$, exactly as in the Burgers' case \eqref{eq:lax-burgers}.
\end{proposition}

\begin{proof}[Proof sketch]
The proofs of Theorem~\ref{thm:rankine-hugoniot} and
Proposition~\ref{prop:entropy-selects-B} used only integration by parts
and the definition of the weak formulation, independent of the specific
form of $f$, so they carry over verbatim.
\end{proof}

\begin{proposition}[General Riemann solution]
\label{prop:general-riemann}
For $f$ strictly convex, the entropy solution of $u_t+f(u)_x=0$ with data
\eqref{eq:general-riemann-data} is:
\begin{enumerate}
  \item[\textup{(a)}] If $u_\ell>u_r$: the shock \eqref{eq:riemann-shock}
  with speed $s$ given by \eqref{eq:general-rh}.
  \item[\textup{(b)}] If $u_\ell \leq u_r$: the rarefaction fan
  \begin{equation}
    \label{eq:general-rarefaction}
    u(x,t) = \begin{cases}
      u_\ell, & x \le f'(u_\ell)\, t, \\
      (f')^{-1}(x/t), & f'(u_\ell)\, t < x < f'(u_r)\, t, \\
      u_r, & x \ge f'(u_r)\, t,
    \end{cases}
  \end{equation}
  where $(f')^{-1}$ is well defined since $f'$ is strictly increasing.
\end{enumerate}
\end{proposition}

\begin{proof}[Proof sketch]
The self-similar ansatz $u(x,t)=g(x/t)$ substituted into $u_t+f(u)_x=0$
gives $g'(\xi)\big(f'(g(\xi))-\xi\big)=0$ (the general-flux analogue of
the computation in Section~\ref{subsec:rarefaction-general}), so either
$g$ is constant or $f'(g(\xi))=\xi$, i.e. $g(\xi)=(f')^{-1}(\xi)$. Taking
$\xi = f'(u_\ell)$ and $\xi=f'(u_r)$ as the two matching boundaries
recovers \eqref{eq:general-rarefaction}; continuity at both boundaries
and vanishing of the corresponding boundary terms in the weak formulation
follow exactly as in the proof of Theorem~\ref{thm:riemann-solution}(b).
Part (a) follows exactly as in Theorem~\ref{thm:riemann-solution}(a),
using Proposition~\ref{prop:general-rh-lax} in place of
\eqref{eq:lax-burgers}.
\end{proof}

\begin{remark}
Taking $f(u)=u^2/2$ gives $f'(u)=u$ and $(f')^{-1}(\xi)=\xi$, recovering
Theorem~\ref{thm:riemann-solution} exactly. The role of convexity is
essential: it is what guarantees $f'$ is strictly increasing, which is in
turn what makes the Lax condition equivalent to the simple inequality
$u_\ell>u_r$, and what makes $(f')^{-1}$ well defined in
\eqref{eq:general-rarefaction}. For non-convex $f$, both the entropy
condition and the Riemann solution are more intricate (they generally
involve the convex hull of $f$ between $u_\ell$ and $u_r$), and are not
treated here.
\end{remark}

\subsection{A numerical scheme: Godunov's method}
\label{subsec:godunov}

This extension outlines a numerical scheme for the conservation law
$u_t+f(u)_x=0$ that uses the exact Riemann solution of
Theorem~\ref{thm:riemann-solution} as its building block, and states its
convergence to the entropy solution.

\subsubsection{Finite volume discretisation}
\label{subsubsec:finite-volume}

Divide $\mathbb{R}$ into cells $[x_{j-1/2}, x_{j+1/2}]$ of width $\Delta
x$, with cell averages
\begin{equation}
  U_j^n \approx \frac{1}{\Delta x}\int_{x_{j-1/2}}^{x_{j+1/2}} u(x, t^n)\, dx,
\end{equation}
at times $t^n = n\Delta t$. Integrating $u_t+f(u)_x=0$ over a cell and
over $[t^n,t^{n+1}]$ gives the exact identity
\begin{equation}
  U_j^{n+1} = U_j^n - \frac{\Delta t}{\Delta x}
  \left( \widehat{f}_{j+1/2} - \widehat{f}_{j-1/2} \right),
\end{equation}
where $\widehat f_{j+1/2}$ is the (time-averaged) flux through the
interface $x_{j+1/2}$, which must be approximated by a \emph{numerical
flux function} $\widehat f_{j+1/2} = \widehat f(U_j^n, U_{j+1}^n)$.

\subsubsection{The Godunov flux}
\label{subsubsec:godunov-flux}

Godunov's method defines $\widehat f$ by solving, at each interface, the
Riemann problem with left and right states $U_j^n$ and $U_{j+1}^n$
\emph{exactly}, using Theorem~\ref{thm:riemann-solution}, and evaluating
the resulting self-similar solution along $x/t=0$:
\begin{equation}
  \label{eq:godunov-flux}
  \widehat f_{j+1/2} = f\big(u^R(0;\, U_j^n, U_{j+1}^n)\big),
\end{equation}
where $u^R(\xi; u_\ell, u_r)$ denotes the value at self-similar
coordinate $\xi=x/t$ of the entropy solution from
Theorem~\ref{thm:riemann-solution} with data $(u_\ell,u_r)$. Concretely,
by \eqref{eq:riemann-shock}--\eqref{eq:riemann-rarefaction},
\begin{equation}
  u^R(0; u_\ell,u_r) =
  \begin{cases}
    u_\ell, & u_\ell>u_r,\ s=\tfrac{u_\ell+u_r}{2}>0, \\
    u_r, & u_\ell>u_r,\ s<0, \\
    u_\ell, & u_\ell\le u_r,\ u_\ell\ge0, \\
    u_r, & u_\ell\le u_r,\ u_r\le0, \\
    0, & u_\ell < 0 < u_r,
  \end{cases}
\end{equation}
the last case corresponding to $\xi=0$ lying inside the rarefaction fan of
\eqref{eq:riemann-rarefaction}, i.e. $u^R(0)=(f')^{-1}(0)=0$ for Burgers'
flux.

\subsubsection{Stability and convergence}
\label{subsubsec:godunov-convergence}

\begin{proposition}[CFL condition]
\label{prop:cfl}
The Godunov scheme is stable provided the time step satisfies
\begin{equation}
  \label{eq:cfl}
  \Delta t \le \frac{\Delta x}{\max_j |f'(U_j^n)|}.
\end{equation}
\end{proposition}

Condition \eqref{eq:cfl} (the Courant--Friedrichs--Lewy condition)
ensures that information does not propagate across more than one cell
per time step, consistent with the finite propagation speed of
characteristics established in Section~\ref{sec:characteristics}.

\begin{theorem}[Convergence; see \cite{evans2010} and standard finite
volume references]
\label{thm:godunov-convergence}
Under \eqref{eq:cfl}, the Godunov approximations converge, as $\Delta x,
\Delta t \to 0$, to the unique entropy solution of
Theorem~\ref{thm:oleinik-uniqueness}.
\end{theorem}

\begin{remark}
This gives a second, independent route to the solutions constructed by
hand in Sections~\ref{sec:nonuniqueness}--\ref{sec:riemann-problem}: an
implementation of \eqref{eq:godunov-flux} applied to the Riemann data
\eqref{eq:riemann-data} should reproduce the shock speed $s=1/2$ and the
linear profile $u=x/t$ inside the rarefaction fan, providing a
computational check on the exact formulas derived analytically above.
\end{remark}

\subsubsection{Numerical results}
\label{subsubsec:godunov-results}

The scheme of Section~\ref{subsubsec:finite-volume}--\ref{subsubsec:godunov-flux}
was implemented with $N=400$ cells on $x\in[-2,2]$, using the CFL number
$0.45$ in \eqref{eq:cfl}, and run to $t=1$ on both Riemann problems
considered earlier: the rarefaction data $(u_\ell,u_r)=(0,1)$ of
Section~\ref{sec:nonuniqueness}, and the shock data $(u_\ell,u_r)=(1,0)$.
Figure~\ref{fig:godunov-results} compares the numerical solution against
the exact entropy solution of Theorem~\ref{thm:riemann-solution}.

\begin{figure}[h]
\centering
\includegraphics[width=0.48\textwidth]{figures/godunov_rarefaction.pdf}
\includegraphics[width=0.48\textwidth]{figures/godunov_shock.pdf}
\caption{Godunov scheme (blue dots, $N=400$ cells) against the exact
entropy solution (red) of Theorem~\ref{thm:riemann-solution}, at
$t=1$. Left: rarefaction data $(u_\ell,u_r)=(0,1)$; the scheme
reproduces the linear fan $u=x/t$, with mild smearing at the two corners
typical of a first-order scheme. Right: shock data $(u_\ell,u_r)=(1,0)$;
the discontinuity is captured sharply at $x=st=0.5$, exactly the
Rankine--Hugoniot speed from \eqref{eq:rankine-hugoniot-burgers}, with
no spurious oscillation.}
\label{fig:godunov-results}
\end{figure}

The measured $L^1$ errors between the numerical and exact solutions are
approximately $4.5\times10^{-3}$ (rarefaction case) and $1.2\times10^{-3}$
(shock case), consistent with the first-order convergence rate of
Theorem~\ref{thm:godunov-convergence}; the shock case converges faster
per cell because the error is concentrated in the single cell straddling
the discontinuity, whereas the rarefaction error accumulates gradually
across the two fan corners.

\subsection{Interaction of two shocks}
\label{subsec:shock-interaction}

This extension studies what happens when two shocks, moving at different
speeds, collide. Remarkably, the answer follows directly from
Theorem~\ref{thm:riemann-solution} without any new machinery.

\subsubsection{Setup}
\label{subsubsec:interaction-setup}

Consider three constant states $u_1 > u_2 > u_3$ and initial data
\begin{equation}
  \label{eq:three-state-data}
  u_0(x) = \begin{cases} u_1, & x<a, \\ u_2, & a<x<b, \\ u_3, & x>b, \end{cases}
  \qquad a<b.
\end{equation}
Since $u_1>u_2$ and $u_2>u_3$, both jumps individually satisfy the Lax
condition \eqref{eq:lax-burgers}, so for small $t>0$ the entropy solution
consists of two shocks, moving at the Rankine--Hugoniot speeds
\begin{equation}
  \label{eq:two-shock-speeds}
  s_{12} = \frac{u_1+u_2}{2}, \qquad s_{23} = \frac{u_2+u_3}{2},
\end{equation}
with positions $x_L(t) = a+s_{12}\,t$ and $x_R(t) = b+s_{23}\,t$, exactly
as in Theorem~\ref{thm:riemann-solution}(a) applied locally at each jump.

\subsubsection{Collision time}
\label{subsubsec:collision-time}

\begin{proposition}[Shock collision]
\label{prop:shock-collision}
Since $u_1>u_2>u_3$, $s_{12}>s_{23}$, so the two shocks collide at the
finite time
\begin{equation}
  \label{eq:collision-time}
  t^{\times} = \frac{b-a}{s_{12}-s_{23}} = \frac{2(b-a)}{u_1-u_3},
\end{equation}
at the point $x^{\times} = a + s_{12}\, t^{\times}$.
\end{proposition}

\begin{proof}
$s_{12}-s_{23} = \tfrac12(u_1+u_2) - \tfrac12(u_2+u_3) = \tfrac12(u_1-u_3)
> 0$, so $x_L(t)$ and $x_R(t)$ are both affine in $t$ with $x_L$
increasing faster than $x_R$ and $x_L(0)=a<b=x_R(0)$; they therefore meet
at a unique finite time, found by solving $x_L(t^\times)=x_R(t^\times)$,
which gives \eqref{eq:collision-time}.
\end{proof}

\subsubsection{The solution after collision}
\label{subsubsec:post-collision}

\begin{theorem}[Merged shock]
\label{thm:merged-shock}
For $t \geq t^\times$, the entropy solution of \eqref{eq:three-state-data}
is the single shock
\begin{equation}
  \label{eq:merged-shock}
  u(x,t) = \begin{cases} u_1, & x < x^\times + s_{13}(t-t^\times), \\
                          u_3, & x > x^\times + s_{13}(t-t^\times), \end{cases}
  \qquad s_{13} = \frac{u_1+u_3}{2}.
\end{equation}
\end{theorem}

\begin{proof}
At $t=t^\times$, the middle state $u_2$ occupies the shrinking interval
$[x_L(t), x_R(t)]$, which by construction has width $0$ exactly at
$t^\times$; hence at $t=t^\times$ the solution is precisely the two-state
Riemann data $u(x,t^\times) = u_1$ for $x<x^\times$, $u_3$ for
$x>x^\times$, with no intermediate state remaining. Since $u_1>u_3$, this
is exactly the hypothesis of Theorem~\ref{thm:riemann-solution}(a)
(shifted to start at time $t^\times$ and position $x^\times$ rather than
$0$ and $0$, which does not affect the argument, as the equation is
autonomous in both $x$ and $t$), whose unique entropy solution for
$t\geq t^\times$ is exactly the shock \eqref{eq:merged-shock}.
\end{proof}

\begin{remark}
Theorem~\ref{thm:merged-shock} illustrates a general principle: an
entropy solution's dependence on only the two bounding Riemann states
means any intermediate structure that shrinks to a point is simply
forgotten. This is also why the Godunov scheme of
Section~\ref{subsec:godunov} handles merging shocks with no special
logic: each time step re-solves a fresh Riemann problem at every
interface.
\end{remark}

\subsection{A damped Burgers' equation}
\label{subsec:damped-burgers}

This extension adds a linear damping term,
\begin{equation}
  \label{eq:damped-burgers}
  u_t + u\, u_x = -\beta\, u, \qquad \beta > 0,
\end{equation}
and asks how damping affects the shock formation result of
Proposition~\ref{prop:shock-time}. Equation \eqref{eq:damped-burgers}
arises, for instance, as a model for weakly nonlinear waves subject to
linear dissipation.

\subsubsection{Characteristics with damping}
\label{subsubsec:damped-characteristics}

Along a curve $x(t)$ with $\dot x(t) = u(x(t),t)$, the chain rule together
with \eqref{eq:damped-burgers} gives
\begin{equation}
  \frac{d}{dt}\, u(x(t),t) = u_t + u\, u_x = -\beta\, u(x(t),t),
\end{equation}
so $u$ no longer stays constant along characteristics, but decays
exponentially:
\begin{equation}
  \label{eq:damped-u-along-char}
  u(x(t),t) = u_0(x_0)\, e^{-\beta t}.
\end{equation}
Substituting \eqref{eq:damped-u-along-char} into $\dot x(t)=u(x(t),t)$ and
integrating gives the characteristic curves explicitly:
\begin{equation}
  \label{eq:damped-characteristic}
  x(t) = x_0 + u_0(x_0)\, \frac{1-e^{-\beta t}}{\beta}.
\end{equation}
Unlike the undamped case, these are no longer straight lines, but they
remain explicit, and the implicit solution formula generalises exactly
as in Theorem~\ref{thm:implicit-solution}: $u(x,t) = u_0(x_0)e^{-\beta t}$
where $x$ and $x_0$ are related by \eqref{eq:damped-characteristic},
valid wherever
\begin{equation}
  \label{eq:damped-phi}
  \Phi_t^\beta(x_0) := x_0 + u_0(x_0)\,\frac{1-e^{-\beta t}}{\beta}
\end{equation}
is invertible.

\subsubsection{A threshold for shock formation}
\label{subsubsec:damping-threshold}

\begin{theorem}[Damping can prevent shock formation]
\label{thm:damped-shock-threshold}
Let $m = \min_{x_0} u_0'(x_0)$. If $m \geq -\beta$, the classical solution
of \eqref{eq:damped-burgers} exists for all $t \geq 0$. If $m < -\beta$,
a shock forms at the finite time
\begin{equation}
  \label{eq:damped-shock-time}
  T^*_\beta = -\frac{1}{\beta} \ln\!\left(1 + \frac{\beta}{m}\right).
\end{equation}
\end{theorem}

\begin{proof}
As in the proof of Proposition~\ref{prop:shock-time}, breakdown occurs at
the first time some $(\Phi_t^\beta)'(x_0) = 1 + u_0'(x_0)\,
\frac{1-e^{-\beta t}}{\beta}$ vanishes. Writing $\tau(t) =
(1-e^{-\beta t})/\beta$, this requires $\tau(t) = -1/u_0'(x_0)$ for some
$x_0$, which (since $\tau$ is increasing) is hardest to achieve for the
most negative value of $u_0'$, namely $x_0$ attaining $m$. Since $\tau$
increases monotonically from $\tau(0)=0$ towards $\tau(\infty) =
1/\beta$ without ever reaching it, a finite solution $t$ to $\tau(t) =
-1/m$ exists if and only if $-1/m < 1/\beta$, i.e. $m < -\beta$; solving
$\tau(T_\beta^*) = -1/m$ for $T_\beta^*$ gives
\eqref{eq:damped-shock-time}. If instead $m \geq -\beta$, then $-1/m \geq
1/\beta \geq \tau(t)$ for every finite $t$ (interpreting $-1/m=\infty$ if
$m=0$), so $(\Phi_t^\beta)'(x_0) > 0$ for every $x_0$ and every $t\geq0$,
and the classical solution persists for all time by the same argument as
Theorem~\ref{thm:implicit-solution}.
\end{proof}

\begin{remark}
Setting $\beta=0$ recovers Proposition~\ref{prop:shock-time} exactly, as
can be checked with L'H\^opital's rule on $-\frac1\beta\ln(1+\beta/m)$
as $\beta\to0$. More strikingly, the theorem shows that for \emph{any}
initial data, however steep, taking $\beta\geq -m$ prevents shock
formation entirely: sufficiently strong linear damping smooths out the
steepening mechanism of Section~\ref{sec:shock-formation} before
characteristics can converge.
\end{remark}

\subsection{Long-time behaviour: the N-wave}
\label{subsec:n-wave}

This final extension describes the universal shape that entropy
solutions with compactly supported initial data approach as $t \to
\infty$, regardless of the detailed shape of $u_0$.

\subsubsection{An exact self-similar solution}
\label{subsubsec:n-wave-exact}

\begin{proposition}[The N-wave]
\label{prop:n-wave-exact}
Fix constants $P, Q > 0$. The function
\begin{equation}
  \label{eq:n-wave}
  N(x,t) = \begin{cases}
    0, & x \le -\sqrt{2Qt}, \\
    x/t, & -\sqrt{2Qt} < x < \sqrt{2Pt}, \\
    0, & x \ge \sqrt{2Pt},
  \end{cases}
  \qquad t>0,
\end{equation}
is an entropy solution of $u_t+u\,u_x=0$ on $\mathbb{R}\times(0,\infty)$.
\end{proposition}

\begin{proof}
On the open middle interval, $N(x,t)=x/t$ is exactly the self-similar
rarefaction profile of Section~\ref{subsec:rarefaction-general}, which
solves the PDE classically there; the outer regions are trivially
constant solutions. At the right edge $x_+(t)=\sqrt{2Pt}$, the one-sided
limits are $u_\ell = x_+(t)/t = \sqrt{2P/t}$ and $u_r=0$; since
$\dot x_+(t) = \sqrt{P/(2t)} = \tfrac12\big(\sqrt{2P/t}+0\big)$, this is
exactly the Rankine--Hugoniot speed \eqref{eq:rankine-hugoniot-burgers},
and $u_\ell > u_r$ so the Lax condition \eqref{eq:lax-burgers} holds. The
same check at the left edge $x_-(t)=-\sqrt{2Qt}$, where $u_\ell=0$ and
$u_r=-\sqrt{2Q/t}$, gives $\dot x_-(t) = -\sqrt{Q/(2t)} =
\tfrac12\big(0+(-\sqrt{2Q/t})\big)$, again matching
\eqref{eq:rankine-hugoniot-burgers}, with $u_\ell=0>-\sqrt{2Q/t}=u_r$
satisfying the Lax condition. Hence both jumps are entropy-admissible,
and $N$ is an entropy solution by Theorem~\ref{thm:rankine-hugoniot} and
Definition~\ref{def:lax-entropy}.
\end{proof}

The profile \eqref{eq:n-wave} is called an \emph{N-wave}: it rises
linearly through the origin and is bounded by a shock on each side, so
that its graph traces the shape of the letter N (reversed), decaying in
amplitude like $t^{-1/2}$ and spreading in width like $t^{1/2}$ as
$t\to\infty$. It is an exact solution for every $t>0$, though its formal
limit as $t\to0^+$ is not an ordinary function of the type considered in
Section~\ref{sec:weak-solutions}: both edges collapse to the origin,
consistent with $N(\cdot,t)$ representing, in the limit, a state
concentrated entirely at a point.

\subsubsection{Convergence of general solutions to the N-wave}
\label{subsubsec:n-wave-convergence}

\begin{theorem}[Asymptotic N-wave decay; Lax, 1957 \cite{lax1957}]
\label{thm:n-wave-decay}
Let $u_0 \in L^1(\mathbb{R})$ have compact support, and set
\begin{equation}
  \label{eq:pq-areas}
  P = \int_{\mathbb{R}} \max(u_0(x),0)\, dx, \qquad
  Q = \int_{\mathbb{R}} \max(-u_0(x),0)\, dx.
\end{equation}
Then the entropy solution $u$ of \eqref{eq:ivp} satisfies
\begin{equation}
  \|u(\cdot,t) - N(\cdot,t)\|_{L^1(\mathbb{R})} \to 0 \qquad \text{as } t \to \infty,
\end{equation}
where $N$ is the N-wave \eqref{eq:n-wave} with the same constants $P,Q$.
\end{theorem}

The proof is not reproduced here, but the mechanism is exactly the
shock-merging phenomenon of Section~\ref{subsec:shock-interaction}: a
compactly supported $u_0$ produces, for small $t$, a solution with
possibly many shocks and rarefactions, but by
Theorem~\ref{thm:merged-shock} every pair of shocks that meet coalesces
into one, so only finitely many discontinuities survive as $t\to\infty$.
Since the total signed area $\int u(x,t)\,dx = P-Q$ is conserved for all
$t$ (integrate \eqref{eq:weak-formulation} over all $x$), the surviving
structure is forced into the unique self-similar shape consistent with
these boundary areas --- exactly \eqref{eq:n-wave}.

\begin{remark}
This result explains the term ``N-wave'' as used in acoustics for the
pressure signature of a sonic boom: an initially localised disturbance,
regardless of its detailed shape, is predicted by the same mechanism to
decay into the universal N-shaped profile of
Proposition~\ref{prop:n-wave-exact} once nonlinear and dissipative
effects have had enough time to act, with amplitude decaying as
$t^{-1/2}$, exactly as observed for sonic booms far from their source.
\end{remark}